% !TeX program = xelatex
\documentclass[11pt,a4paper]{article}

\usepackage[a4paper,margin=2.2cm]{geometry}
\usepackage{amsmath,amssymb,bm}
\usepackage{booktabs}
\usepackage{array}
\usepackage{hyperref}
\usepackage{xcolor}
\usepackage{graphicx}

\hypersetup{
  colorlinks=true,
  linkcolor=blue,
  urlcolor=blue,
  citecolor=blue
}

\newcommand{\Tr}{\operatorname{Tr}}

\title{\textbf{Deuteron Spin Density Matrix and the $n$--$p$ Spatial Distribution}\\
\large A Vertical Pure-Tensor Deuteron Beam at SAMURAI}
\author{}
\date{}

\begin{document}
\maketitle

\begin{abstract}
The deuteron $S$--$D$ coupling links spin polarization to the angular
distribution of the proton--neutron relative coordinate.
For a beam axially symmetric about the vertical $y$ axis, the spatial
density has a quadrupole term proportional to
$p_{yy}[\sqrt2\,uw-w^2/2]P_2(\cos\theta_y)$, where $u,w$ are radial wave functions.
Where this radial coefficient is positive, $p_{yy}=+1$ enhances the
probability along $y$, while $p_{yy}=-2$ enhances it in the perpendicular
plane. Both modes have the same radial separation distribution.
This relation follows from the beam spin density matrix and the deuteron
wave function and provides a coordinate-sampling prescription for
semiclassical models.
\end{abstract}

\section{Spin--orbital coupling in the deuteron wave function}

The deuteron ground state has $J^\pi=1^+$.
In the nonrelativistic description, it is dominated by the ${}^3S_1$
component, which the tensor part of the nuclear interaction couples to
${}^3D_1$. For a chosen quantization axis, the state with
total-angular-momentum projection $M$ is
\begin{equation}
\Psi_{1M}(\bm r,\sigma)
=
\frac{u(r)}{r}Y_{00}(\hat{\bm r})\chi_{1M}(\sigma)
+
\frac{w(r)}{r}
\sum_{m_L,m_S}
\langle 2m_L,1m_S|1M\rangle
Y_{2m_L}(\hat{\bm r})\chi_{1m_S}(\sigma),
\label{eq:deuteronwf}
\end{equation}
where $\bm r=\bm r_p-\bm r_n$ is the proton--neutron relative coordinate,
$\sigma$ denotes the nucleon-spin degrees of freedom, and $\chi_{1m_S}$
is an $S=1$ spin-triplet state. The $S$- and $D$-wave radial functions
$u(r)$ and $w(r)$ satisfy
\[
\int_0^\infty [u^2(r)+w^2(r)]\,dr=1.
\]

In the pure $S$-wave limit, $w=0$, the orbital factor $Y_{00}$ is
spherically symmetric and independent of $M$. Any magnetic-substate
populations therefore give a spherical spatial distribution.
The $D$ wave correlates $m_L$ and $m_S$ through Clebsch--Gordan
coefficients, so the total wave function generally cannot factorize
into separate orbital and spin parts. This coupling gives different
magnetic substates different spatial angular distributions.

\section{Magnetic-substate populations and beam polarization}

Take the vertical $y$ direction as the quantization axis and use the basis
$\{|+1\rangle_y,|0\rangle_y,|-1\rangle_y\}$.
The beam's $3\times3$ spin density matrix describes populations and
coherences among these three $J=1$ magnetic substates:
\begin{equation}
\rho^{(J=1)}
=
\begin{pmatrix}
\rho_{++} & \rho_{+0} & \rho_{+-}\\
\rho_{0+} & \rho_{00} & \rho_{0-}\\
\rho_{-+} & \rho_{-0} & \rho_{--}
\end{pmatrix},
\qquad
\Tr\rho^{(J=1)}=1.
\label{eq:spinrho}
\end{equation}
Axial symmetry about $y$ makes the matrix diagonal in this basis:
\begin{equation}
\rho^{(J=1)}=\operatorname{diag}(n_+,n_0,n_-),
\qquad n_++n_0+n_-=1.
\end{equation}
Define the vector and tensor polarizations as
\begin{align}
p_y &= n_+-n_-,\\
p_{yy} &= n_++n_--2n_0,
\label{eq:polarization}
\end{align}
with ranges $-1\le p_y\le1$ and $-2\le p_{yy}\le1$.
Pure-tensor polarization satisfies $p_y=0$ and $p_{yy}\ne0$.

\section{Tracing over nucleon spin to obtain the spatial density}

Combining the beam density matrix with the complete wave functions in
Eq.~\eqref{eq:deuteronwf} gives the statistical operator for the orbital
and nucleon-spin degrees of freedom:
\begin{equation}
\hat\rho_{\rm tot}
=
\sum_{M,M'=-1}^{1}
\rho_{MM'}^{(J=1)}|\Psi_M\rangle\langle\Psi_{M'}|.
\label{eq:fullrho}
\end{equation}
Tracing over the unobserved nucleon spins gives the reduced spatial
density matrix,
\begin{equation}
\rho_{\rm sp}(\bm r,\bm r')
=
\sum_\sigma
\langle\bm r,\sigma|\hat\rho_{\rm tot}|\bm r',\sigma\rangle.
\label{eq:reduced}
\end{equation}
Its diagonal element is the probability density per unit volume:
\begin{equation}
\varrho(\bm r)
=
\rho_{\rm sp}(\bm r,\bm r)
=
\sum_{M,M'}\rho_{MM'}^{(J=1)}
\sum_\sigma\Psi_M(\bm r,\sigma)\Psi_{M'}^*(\bm r,\sigma).
\label{eq:localdensity}
\end{equation}
The spatial distribution thus depends jointly on substate populations,
coherences, and the internal deuteron wave function.
The reduced density matrix also describes mixed beams; even a pure
complete deuteron state can yield a mixed spatial state after the trace
because orbital and spin degrees of freedom are entangled.

\section{Spatial distribution of an axisymmetric beam}

Define
\[
\cos\theta_y=\hat{\bm r}\cdot\hat{\bm y},
\qquad P_2(\cos\theta_y)=\frac12(3\cos^2\theta_y-1),
\]
and the radial quadrupole coefficient
\begin{equation}
A(r)=\sqrt2\,u(r)w(r)-\frac12w^2(r).
\label{eq:A}
\end{equation}
Substituting Eq.~\eqref{eq:deuteronwf} into Eq.~\eqref{eq:localdensity}
and summing over nucleon spin gives
\begin{align}
\varrho_0(r,\theta_y)
&=\frac{u^2+w^2-2A(r)P_2(\cos\theta_y)}{4\pi r^2},
\label{eq:rho0}\\
\varrho_{\pm1}(r,\theta_y)
&=\frac{u^2+w^2+A(r)P_2(\cos\theta_y)}{4\pi r^2}.
\label{eq:rhopm}
\end{align}
Weighting by the populations,
$\varrho=n_+\varrho_{+1}+n_0\varrho_0+n_-\varrho_{-1}$, yields
\begin{equation}
\boxed{
\varrho(r,\theta_y)
=\frac{1}{4\pi r^2}
\left[u^2(r)+w^2(r)+p_{yy}A(r)P_2(\cos\theta_y)\right].
}
\label{eq:master}
\end{equation}
Since $\varrho_{+1}=\varrho_{-1}$, the weighted sum depends only on
$n_++n_-$. The vector polarization $p_y=n_+-n_-$ therefore drops out of
the scalar coordinate density. Spin-dependent scattering can still be
sensitive to vector polarization.

\subsection{\texorpdfstring{$S$--$D$}{S--D} interference sets the quadrupole strength}

The coefficient $A(r)$ contains both the interference term
$\sqrt2\,uw$ and the $D$-wave term $-w^2/2$.
Where $|w|\ll|u|$, the interference is first order in $w$, so a small
$D$-state probability can still produce an appreciable angular variation.
The $S$ and $D$ components in Eq.~\eqref{eq:deuteronwf} form a coherent
superposition, and their interference must be retained in the spatial density.

\subsection{Radial distribution and the unpolarized limit}

Using $\int P_2(\cos\theta_y)\,d\Omega=0$, angular integration of
Eq.~\eqref{eq:master} gives
\begin{equation}
r^2\int d\Omega\,\varrho(r,\Omega)=u^2(r)+w^2(r).
\label{eq:radial}
\end{equation}
For fixed bound-state wave functions $u,w$, the radial probability is
$P(r)\,dr=[u^2(r)+w^2(r)]\,dr$ for all substate populations.
Polarization changes the directional probability at a given $r$ through
the quadrupole term; the radial separation distribution remains unchanged.

An unpolarized beam has $n_+=n_0=n_-=1/3$, hence $p_y=p_{yy}=0$.
The quadrupole terms of the three substates cancel, leaving the spherical
density $\varrho=(u^2+w^2)/(4\pi r^2)$.

\section{Two ideal pure-tensor modes}

Both modes have $p_y=0$. Their populations and spatial distributions
are listed below; the stated directional enhancements assume $A(r)>0$.
\begin{center}
\begin{tabular}{>{\centering\arraybackslash}p{2.3cm}
                >{\centering\arraybackslash}p{3.4cm}
                p{7.1cm}}
\toprule
Beam mode & $(n_+,n_0,n_-)$ & Spatial distribution ($A(r)>0$)\\
\midrule
$p_{yy}=+1$ & $(1/2,0,1/2)$ & Enhanced along $y$; prolate\\
$p_{yy}=-2$ & $(0,1,0)$ & Enhanced in the $x$--$z$ plane; oblate\\
Unpolarized & $(1/3,1/3,1/3)$ & Quadrupole terms cancel; spherical\\
\bottomrule
\end{tabular}
\end{center}

The $p_{yy}=+1$ mode is an equal incoherent mixture of $M_y=+1$ and
$M_y=-1$. The $p_{yy}=-2$ mode is the pure state $|J=1,M_y=0\rangle$.
The term ``pure tensor'' specifies the polarization components;
quantum-state purity is determined by the density matrix.

Along $y$, $P_2(1)=1$; in the $x$--$z$ plane, $P_2(0)=-1/2$.
Substituting these values into Eq.~\eqref{eq:master} gives the directional
enhancements in the table. Wherever $A(r)<0$, these enhancements reverse.
Prolate and oblate describe angular anisotropy at a given $r$;
both polarization modes have the same mean $n$--$p$ separation.

\clearpage
\section{Probability-density comparison of the two modes}

Figure~\ref{fig:rnp-density} compares the relative-coordinate density for
$p_{yy}=+1$ and $-2$: the spatial anisotropy changes, while the radial
curves coincide. To make Eq.~\eqref{eq:master} concrete without selecting
a nuclear potential, use the normalized illustrative model
\begin{equation}
u(r)=\frac{2\sqrt{1-P_D}}{\pi^{1/4}b^{3/2}}r e^{-r^2/(2b^2)},
\qquad
w(r)=\frac{4\sqrt{P_D}}{\sqrt{15}\,\pi^{1/4}b^{7/2}}r^3 e^{-r^2/(2b^2)},
\label{eq:plot-model}
\end{equation}
with $b=2\,\mathrm{fm}$ and $P_D=\int w^2\,dr=0.05$.
These are chosen illustration parameters, not a fitted deuteron wave function.
Writing $r_{np}=|\bm r_p-\bm r_n|=r$ and $\mu=\cos\theta_y$, the
angle marginal integrated over radius and azimuth is
\begin{equation}
\frac{dP}{d\mu}=\frac12[1+p_{yy}C P_2(\mu)],\qquad
C=\int_0^\infty A(r)\,dr
=\sqrt{\frac65P_D(1-P_D)}-\frac{P_D}{2}.
\end{equation}
The program \texttt{plots/plot\_rnp\_density.py} generates the figure;
running it with \texttt{python3} from this directory uses these defaults.

\begin{figure}[htbp]
\centering
\includegraphics[width=0.94\textwidth]{plots/output/rnp_density_pyy.pdf}
\caption{Illustrative $n$--$p$ probability densities.
(a,b) $z_{np}=0$ slices of the three-dimensional density, with a common
color scale in $\mathrm{fm}^{-3}$; these are not projections integrated over $z$.
(c) The radial density $P(r_{np})=u^2+w^2$ is identical for both modes.
(d) The normalized density in $\mu=\cos\theta_y$ distinguishes alignment
along $y$ ($\mu=\pm1$) from the perpendicular plane ($\mu=0$).}
\label{fig:rnp-density}
\end{figure}

\clearpage
\subsection{Three-dimensional constant-density surfaces}

At the same absolute density, the $p_{yy}=+1$ surface extends farther
along $y$, whereas the $p_{yy}=-2$ surface extends farther in the
$x$--$z$ plane (Fig.~\ref{fig:rnp-isosurfaces}).
The surfaces solve $\varrho(r,\theta_y)=0.003\,\mathrm{fm}^{-3}$ using
Eq.~\eqref{eq:plot-model}; they retain the common spatial scale and
vertical $y$ axis. The transparent overlay compares their geometry,
not a new mixed-state density or a surface enclosing a prescribed probability.
The program \texttt{plots/plot\_rnp\_isosurfaces.py} generates this comparison.

\begin{figure}[htbp]
\centering
\includegraphics[width=\textwidth]{plots/output/rnp_isosurfaces_pyy.pdf}
\caption{Constant-density surfaces of the illustrative S--D model at
$\varrho=0.003\,\mathrm{fm}^{-3}$.
(a) $p_{yy}=+1$, blue; (b) $p_{yy}=-2$, orange;
(c) both surfaces overlaid with opacity $0.24$ per surface.
All panels use the same orthographic view and equal axis scales in fm.
The relative coordinate is $\bm r_{np}=\bm r_p-\bm r_n$.}
\label{fig:rnp-isosurfaces}
\end{figure}

\section{Substate coherence and azimuthal dependence}

Off-diagonal elements $\rho_{MM'}$ with $M\ne M'$ describe coherence
between magnetic substates. Their tensor components can produce
azimuthal dependence in $\phi$, which requires the full expression in
Eq.~\eqref{eq:localdensity}. The single $p_{yy}P_2(\cos\theta_y)$ term in
Eq.~\eqref{eq:master} applies to axial symmetry.

A general spin-1 density matrix decomposes into rank-0, rank-1, and
rank-2 components, corresponding to normalization, vector polarization,
and tensor polarization. An axisymmetric pure-tensor beam retains the
normalization component and the tensor component along the chosen axis.

\clearpage
\section{Coordinate sampling for QMD and semiclassical models}

For an axisymmetric beam, relative coordinates can be sampled directly
from Eq.~\eqref{eq:master}. Since $\varrho$ is a density per unit volume,
the probability element in spherical coordinates is
\begin{equation}
 dP=\varrho(r,\theta_y)r^2\,dr\,d\Omega
 =\frac{u^2+w^2+p_{yy}A(r)P_2(\cos\theta_y)}{4\pi}\,dr\,d\Omega,
\label{eq:sampling}
\end{equation}
where $d\Omega=\sin\theta_y\,d\theta_y\,d\phi$.
One can first sample $r$ from $P(r)=u^2+w^2$, then sample the direction from
\[
p(\Omega\mid r)=\frac{1}{4\pi}
\left[1+\frac{p_{yy}A(r)}{u^2(r)+w^2(r)}P_2(\cos\theta_y)\right].
\]
Axial symmetry makes $\phi$ uniform on $[0,2\pi)$.

Fixing every relative coordinate along $y$ or confining it to the
$x$--$z$ plane imposes an artificially strong directional constraint
and can overestimate polarization effects. The sampling prescription
above preserves the coordinate probability of the quantum state.
Reaction models with explicit nucleon-spin dependence also need the
$m_L$--$m_S$ correlations in Eq.~\eqref{eq:deuteronwf}.

\end{document}
